\begin{PSbox}[unbreakable]{obj}{اهداف یادگیری}

پس از تکمیل این ماژول، دانشجویان قادر خواهند بود:

\begin{enumerate}
\def\labelenumi{\arabic{enumi}.}
\tightlist
\item
  مقادیر امید ریاضی را برای متغیرهای تصادفی گسسته و پیوسته محاسبه کنند
\item
  واریانس و انحراف معیار را محاسبه کرده و معنای مهندسی آن‌ها را تفسیر کنند
\item
  ویژگی‌های امید ریاضی و واریانس را در تحلیل مهندسی به کار ببرند
\item
  گشتاورهای مراتب بالاتر را محاسبه و تفسیر کنند
\item
  از تابع مولد گشتاور \lr{(MGF)} استفاده کنند
\item
  گشتاورهای آماری را به کمیت‌های فیزیکی مهندسی مرتبط کنند
\end{enumerate}

\end{PSbox}

\PSsection{4.1}{مقدمه: چرا آماره‌های خلاصه؟}

\PSsubsection{نیاز مهندسی}

یک توزیع احتمال تمام اطلاعات مربوط به یک متغیر تصادفی را در خود دارد. اما برای تصمیم‌های مهندسی، اغلب به خلاصه‌های موجز نیاز داریم:

\begin{itemize}
\tightlist
\item
  \textbf{پردازش سیگنال:} توان میانگین چقدر است؟ \lr{\ensuremath{\to}} میانگین \lr{X\textsuperscript{2}}
\item
  \textbf{کنترل کیفیت:} قطعات چقدر تغییر می‌کنند؟ \lr{\ensuremath{\to}} واریانس
\item
  \textbf{مخابرات:} توان نویز چقدر است؟ \lr{\ensuremath{\to} E[N\textsuperscript{2}]}
\item
  \textbf{قابلیت اطمینان:} عمر میانگین چقدر است؟ \lr{\ensuremath{\to}} میانگین
\end{itemize}

گشتاورها (میانگین، واریانس و \lr{…}) خلاصه‌های ضروری و \textbf{مرتبط با مهندسی} از یک توزیع هستند.

\PSsection{4.2}{مقدار امید ریاضی (میانگین)}

\begin{PSbox}[unbreakable]{def}{تعریف \lr{—} متغیر تصادفی گسسته}

\PSdisplay{E[X] = \mu_X = \sum_{x} x \cdot p_X(x)}

جمع (هر مقدار \ensuremath{\times} احتمال آن).

\end{PSbox}

\begin{PSbox}[unbreakable]{def}{تعریف \lr{—} متغیر تصادفی پیوسته}

\PSdisplay{E[X] = \mu_X = \int_{-\infty}^{\infty} x \cdot f_X(x) \, dx}

\end{PSbox}

\PSsubsection{تفسیر مهندسی}

مقدار امید ریاضی همان \textbf{میانگین درازمدت} است: اگر آزمایش را بارها تکرار کنید و نتایج را میانگین بگیرید، مقدار تقریبی \lr{E[X]} به دست می‌آید.

{\def\LTcaptype{none} % do not increment counter
\begin{longtable}[c]{@{}cc@{}}
\toprule\noalign{}
\rowcolor{PSnavy}\PSth{کمیت مهندسی} & \PSth{معنای مقدار امید ریاضی} \\
\midrule\noalign{}
\endhead
\bottomrule\noalign{}
\endlastfoot
ولتاژ نویز \lr{E[N]} & بایاس/آفست \lr{DC} در نویز \\
دامنه سیگنال \lr{E[A]} & سطح میانگین سیگنال \\
عمر قطعه \lr{E[T]} & میانگین زمان تا خرابی \lr{(MTTF)} \\
تأخیر بسته \lr{E[D]} & تأخیر میانگین \\
خطاهای بیت \lr{E[X]} & تعداد خطای میانگین \\
\end{longtable}
}

\begin{PSbox}[unbreakable]{ex}{مثال: توان میانگین نویز}

برای نویز گاوسی با میانگین صفر \lr{N \textasciitilde{} N(0, \ensuremath{\sigma}\textsuperscript{2})}:

\begin{itemize}
\tightlist
\item
  \lr{E[N] = 0} (بدون مؤلفه \lr{DC})
\item
  توان میانگین = \lr{E[N\textsuperscript{2}] = \ensuremath{\sigma}\textsuperscript{2}} (این همان واریانس است!)
\end{itemize}

\end{PSbox}

\begin{PSbox}[unbreakable]{ex}{مثال: تعداد مورد انتظار ارسال مجدد}

پروتکل \lr{ARQ}: هر انتقال با احتمال \lr{p = 0.8} موفق است.\PSbr{}تعداد تلاش‌ها \lr{X \textasciitilde{} Geometric(0.8)}.\PSbr{}\lr{E[X] = 1/p = 1/0.8 = 1.25} تلاش به‌طور میانگین.

\end{PSbox}

\PSsection{4.3}{مقدار امید ریاضی یک تابع: \lr{E[g(X)]}}

\begin{PSbox}[unbreakable]{thm}{قضیه \lr{LOTUS} (قانون آماردان ناخودآگاه)}

برای یافتن \lr{E[g(X)]}، نیازی به توزیع \lr{Y = g(X)} ندارید:

\textbf{گسسته:}\PSdisplay{E[g(X)] = \sum_{x} g(x) \cdot p_X(x)}

\textbf{پیوسته:}\PSdisplay{E[g(X)] = \int_{-\infty}^{\infty} g(x) \cdot f_X(x) \, dx}

\end{PSbox}

\PSsubsection{کاربردهای مهندسی}

{\def\LTcaptype{none} % do not increment counter
\begin{longtable}[c]{@{}cc@{}}
\toprule\noalign{}
\rowcolor{PSnavy}\PSth{تابع \lr{g(X)}} & \PSth{معنای مهندسی} \\
\midrule\noalign{}
\endhead
\bottomrule\noalign{}
\endlastfoot
\lr{g(X) = X\textsuperscript{2}} & مقدار میانگین مربع (توان) \\
\lr{g(X) = (X - \ensuremath{\mu})\textsuperscript{2}} & واریانس (توان \lr{AC}) \\
\lr{g(X)} = & \lr{X} \\
\lr{g(X) = e\textsuperscript{j\ensuremath{\omega}X}} & تابع مشخصه \\
\end{longtable}
}

\begin{PSbox}[unbreakable]{ex}{مثال: توان میانگین یک سیگنال}

سیگنال \lr{X} دارای \lr{PDF f\textsubscript{X}(x)} است. توان \lr{P = X\textsuperscript{2}}.

\lr{E[P] = E[X\textsuperscript{2}] = \ensuremath{\int} x\textsuperscript{2} f\textsubscript{X}(x) dx}

برای \lr{X \textasciitilde{} N(0, \ensuremath{\sigma}\textsuperscript{2}): E[X\textsuperscript{2}] = \ensuremath{\sigma}\textsuperscript{2}} (توان نویز برابر واریانس است).

برای \lr{X \textasciitilde{} N(\ensuremath{\mu}, \ensuremath{\sigma}\textsuperscript{2}): E[X\textsuperscript{2}] = \ensuremath{\mu}\textsuperscript{2} + \ensuremath{\sigma}\textsuperscript{2}} (توان کل = توان \lr{DC} + توان \lr{AC}).

\end{PSbox}

\PSsection{4.4}{واریانس}

\begin{PSbox}[unbreakable]{def}{تعریف}

\PSdisplay{\PSmtext{Var}(X) = E[(X - \mu_X)^2] = E[X^2] - (E[X])^2}

فرمول دوم (فرمول محاسباتی) اغلب آسان‌تر است.

\end{PSbox}

\PSsubsection{استخراج فرمول محاسباتی}

\lr{Var(X) = E[(X - \ensuremath{\mu})\textsuperscript{2}]}\PSbr{}= \lr{E[X\textsuperscript{2} - 2\ensuremath{\mu}X + \ensuremath{\mu}\textsuperscript{2}]}\PSbr{}= \lr{E[X\textsuperscript{2}] - 2\ensuremath{\mu}E[X] + \ensuremath{\mu}\textsuperscript{2}}\PSbr{}= \lr{E[X\textsuperscript{2}] - 2\ensuremath{\mu}\textsuperscript{2} + \ensuremath{\mu}\textsuperscript{2}}\PSbr{}= \lr{E[X\textsuperscript{2}] - \ensuremath{\mu}\textsuperscript{2}}

\PSsubsection{تفسیر مهندسی}

واریانس \textbf{پراکندگی/پخش‌شدگی/عدم‌قطعیت} را می‌سنجد:

{\def\LTcaptype{none} % do not increment counter
\begin{longtable}[c]{@{}cc@{}}
\toprule\noalign{}
\rowcolor{PSnavy}\PSth{زمینه} & \PSth{معنای واریانس} \\
\midrule\noalign{}
\endhead
\bottomrule\noalign{}
\endlastfoot
سیگنال نویز & توان نویز (برای نویز با میانگین صفر) \\
ساخت و تولید & میزان انحراف قطعات از مقدار نامی \\
اندازه‌گیری‌ها & دقت اندازه‌گیری (کوچک‌تر = دقیق‌تر) \\
سیگنال & مؤلفه توان \lr{AC} \\
عمر & عمر چقدر قابل پیش‌بینی است؟ \\
\end{longtable}
}

\PSsubsection{نکته کلیدی: واریانس = توان در مهندسی}

برای یک سیگنال با میانگین صفر \lr{X}:

\begin{itemize}
\tightlist
\item
  \lr{E[X\textsuperscript{2}] = Var(X) = \ensuremath{\sigma}\textsuperscript{2}} = \textbf{توان کل}
\item
  برای \lr{X \textasciitilde{} N(0, \ensuremath{\sigma}\textsuperscript{2})}: توان نویز = \lr{\ensuremath{\sigma}\textsuperscript{2} = Var(X)}
\end{itemize}

برای یک سیگنال با آفست \lr{DC (\ensuremath{\mu} \ensuremath{\neq} 0)}:

\begin{itemize}
\tightlist
\item
  \lr{E[X\textsuperscript{2}] = \ensuremath{\mu}\textsuperscript{2} + \ensuremath{\sigma}\textsuperscript{2}} = \textbf{توان کل = توان \lr{DC} + توان \lr{AC}}
\item
  \lr{Var(X) = \ensuremath{\sigma}\textsuperscript{2}} = \textbf{فقط توان \lr{AC}}
\end{itemize}

\PSsection{4.5}{انحراف معیار}

\begin{PSbox}[unbreakable]{def}{تعریف}

\PSdisplay{\sigma_X = \sqrt{\PSmtext{Var}(X)}}

\end{PSbox}

\PSsubsection{اهمیت مهندسی}

انحراف معیار \textbf{همان یکای} \lr{X} را دارد و بنابراین مستقیماً قابل تفسیر است:

{\def\LTcaptype{none} % do not increment counter
\begin{longtable}[c]{@{}ccc@{}}
\toprule\noalign{}
\rowcolor{PSnavy}\PSth{متغیر تصادفی} & \PSth{یکای انحراف معیار} & \PSth{تفسیر} \\
\midrule\noalign{}
\endhead
\bottomrule\noalign{}
\endlastfoot
نویز ولتاژ & ولت & ولتاژ نویز مؤثر \lr{(RMS)} \\
جیتر زمان‌بندی & ثانیه & جیتر مؤثر \lr{(RMS)} \\
خطای فاصله & متر & خطای موقعیت مؤثر \lr{(RMS)} \\
دما & \lr{\ensuremath{^{\circ}}C} & نوسان دمای مؤثر \lr{(RMS)} \\
\end{longtable}
}

انحراف معیار برای سیگنال‌های با میانگین صفر همان مقدار \textbf{مؤثر \lr{(RMS)}} است.

\PSsection{4.6}{ویژگی‌های امید ریاضی}

\PSsubsection{خطی‌بودن (بنیادی)}

\PSdisplay{E[aX + b] = aE[X] + b}

\PSdisplay{E[X + Y] = E[X] + E[Y] \quad \PSmtext{(ALWAYS, even if dependent!)}}

\PSdisplay{E\left[\sum_{i=1}^n a_i X_i\right] = \sum_{i=1}^n a_i E[X_i]}

\PSsubsection{یکنوایی}

اگر همواره \lr{X \ensuremath{\ge} 0} باشد، آنگاه \lr{E[X] \ensuremath{\ge} 0}.

اگر همواره \lr{X \ensuremath{\ge} Y} باشد، آنگاه \lr{E[X] \ensuremath{\ge} E[Y]}.

\PSsubsection{برای متغیرهای تصادفی مستقل}

اگر \lr{X} و \lr{Y} مستقل باشند:\PSdisplay{E[XY] = E[X] \cdot E[Y]}

\begin{PSquote}

\lr{\PSwarn{}} این رابطه برای متغیرهای وابسته برقرار نیست!

\end{PSquote}

\begin{PSbox}[unbreakable]{ex}{مثال مهندسی: نویز کل}

یک سیگنال دریافتی: \lr{R = S + N\textsubscript{1} + N\textsubscript{2}} (سیگنال به‌همراه دو منبع نویز)

\lr{E[R] = E[S] + E[N\textsubscript{1}] + E[N\textsubscript{2}] = E[S] + 0 + 0 = E[S]}

(نویز با میانگین صفر، مقدار مورد انتظار سیگنال دریافتی را بایاس نمی‌کند.)

\end{PSbox}

\PSsection{4.7}{ویژگی‌های واریانس}

\PSsubsection{مقیاس‌بندی و انتقال}

\PSdisplay{\PSmtext{Var}(aX + b) = a^2 \PSmtext{Var}(X)}

\begin{itemize}
\tightlist
\item
  افزودن ثابت \lr{b} واریانس را تغییر نمی‌دهد (انتقال، پراکندگی را تغییر نمی‌دهد)
\item
  مقیاس‌بندی با \lr{a}، واریانس را در \lr{a\textsuperscript{2}} ضرب می‌کند
\end{itemize}

\PSsubsection{جمع متغیرهای تصادفی مستقل}

اگر \lr{X\textsubscript{1}, X\textsubscript{2}, …, X\textsubscript{n}} \textbf{مستقل} باشند:

\PSdisplay{\PSmtext{Var}(X_1 + X_2 + \cdots + X_n) = \PSmtext{Var}(X_1) + \PSmtext{Var}(X_2) + \cdots + \PSmtext{Var}(X_n)}

\begin{PSquote}

\lr{\PSwarn{}} واریانس‌ها فقط برای متغیرهای مستقل جمع می‌شوند!

\end{PSquote}

\begin{PSbox}[unbreakable]{ex}{مثال مهندسی: جمع توان نویز}

دو منبع نویز مستقل (\lr{\ensuremath{\sigma}\textsubscript{1}\textsuperscript{2} = 0.01 V\textsuperscript{2}}، \lr{\ensuremath{\sigma}\textsubscript{2}\textsuperscript{2} = 0.04 V\textsuperscript{2}}):

توان نویز کل = \lr{\ensuremath{\sigma}\textsubscript{1}\textsuperscript{2} + \ensuremath{\sigma}\textsubscript{2}\textsuperscript{2} = 0.05 V\textsuperscript{2}}\PSbr{}نویز مؤثر کل = \lr{\ensuremath{\surd}0.05 = 0.224 V}

توجه: مقادیر مؤثر \lr{(RMS)} مستقیماً جمع نمی‌شوند! \lr{(0.1 + 0.2 \ensuremath{\neq} 0.224)}

\end{PSbox}

\PSsubsection{واریانس میانگین نمونه}

برای \lr{n} مشاهده مستقل و هم‌توزیع \lr{(i.i.d.) X\textsubscript{1}, …, X\textsubscript{n}} با واریانس \lr{\ensuremath{\sigma}\textsuperscript{2}}:

\PSdisplay{\PSmtext{Var}(\bar{X}) = \PSmtext{Var}\left(\frac{1}{n}\sum X_i\right) = \frac{\sigma^2}{n}}

اندازه‌گیری‌های بیشتر \lr{\ensuremath{\to}} عدم‌قطعیت کمتر در میانگین.

\PSsection{4.8}{گشتاورهای مراتب بالاتر}

\PSsubsection{گشتاور مرتبه \lr{n}}

\PSdisplay{\mu_n' = E[X^n] = \int_{-\infty}^{\infty} x^n f_X(x) \, dx}

\PSsubsection{گشتاور مرکزی مرتبه \lr{n}}

\PSdisplay{\mu_n = E[(X - \mu)^n]}

\begin{itemize}
\tightlist
\item
  \lr{\ensuremath{\mu}\textsubscript{1} = 0} (همیشه)
\item
  \lr{\ensuremath{\mu}\textsubscript{2} = Var(X)}
\item
  \lr{\ensuremath{\mu}\textsubscript{3}} به \textbf{چولگی} مربوط است
\item
  \lr{\ensuremath{\mu}\textsubscript{4}} به \textbf{کشیدگی} مربوط است
\end{itemize}

\PSsubsection{چولگی (گشتاور مرکزی سوم، نرمال‌شده)}

\PSdisplay{\gamma_1 = \frac{E[(X-\mu)^3]}{\sigma^3}}

{\def\LTcaptype{none} % do not increment counter
\begin{longtable}[c]{@{}ccc@{}}
\toprule\noalign{}
\rowcolor{PSnavy}\PSth{چولگی} & \PSth{شکل} & \PSth{مثال} \\
\midrule\noalign{}
\endhead
\bottomrule\noalign{}
\endlastfoot
\lr{\ensuremath{\gamma}\textsubscript{1} = 0} & متقارن & گاوسی \\
\lr{\ensuremath{\gamma}\textsubscript{1} \textgreater{} 0} & چوله به راست (دنباله بلند در راست) & نمایی \\
\lr{\ensuremath{\gamma}\textsubscript{1} \textless{} 0} & چوله به چپ (دنباله بلند در چپ) & نمایی منفی‌شده \\
\end{longtable}
}

\textbf{اهمیت مهندسی:} چولگی به ما می‌گوید آیا مقادیر حدی در یک جهت محتمل‌تر هستند یا نه. برای قابلیت اطمینان مهم است (توزیع‌های زمان خرابی چوله به راست هستند).

\PSsubsection{کشیدگی (گشتاور مرکزی چهارم، نرمال‌شده)}

\PSdisplay{\gamma_2 = \frac{E[(X-\mu)^4]}{\sigma^4} - 3}

(کسر کردن 3 باعث می‌شود کشیدگی گاوسی صفر شود: «کشیدگی مازاد».)

{\def\LTcaptype{none} % do not increment counter
\begin{longtable}[c]{@{}ccc@{}}
\toprule\noalign{}
\rowcolor{PSnavy}\PSth{کشیدگی} & \PSth{تفسیر} & \PSth{معنا} \\
\midrule\noalign{}
\endhead
\bottomrule\noalign{}
\endlastfoot
\lr{\ensuremath{\gamma}\textsubscript{2} = 0} & مزوکورتیک & دنباله‌هایی شبیه گاوسی \\
\lr{\ensuremath{\gamma}\textsubscript{2} \textgreater{} 0} & لپتوکورتیک & دنباله‌های سنگین‌تر از گاوسی \\
\lr{\ensuremath{\gamma}\textsubscript{2} \textless{} 0} & پلاتی‌کورتیک & دنباله‌های سبک‌تر از گاوسی \\
\end{longtable}
}

\textbf{اهمیت مهندسی:} کشیدگی زیاد به معنای مقادیر پرت بیشتر است \lr{—} برای مدل‌سازی نویز تکانشی مهم است.

\PSsection{4.9}{تابع مولد گشتاور \lr{(MGF)}}

\begin{PSbox}[unbreakable]{def}{تعریف}

\PSdisplay{M_X(t) = E[e^{tX}]}

\textbf{گسسته:} \lr{M\textsubscript{X}(t) = \ensuremath{\Sigma} e\textsuperscript{tx} p\textsubscript{X}(x)}

\textbf{پیوسته:} \lr{M\textsubscript{X}(t) = \ensuremath{\int} e\textsuperscript{tx} f\textsubscript{X}(x) dx}

\end{PSbox}

\PSsubsection{چرا «مولد گشتاور»؟}

گشتاور مرتبه \lr{n} با \lr{n} بار مشتق‌گیری در \lr{t = 0} به دست می‌آید:

\PSdisplay{E[X^n] = \frac{d^n M_X(t)}{dt^n}\bigg|_{t=0}}

\begin{itemize}
\tightlist
\item
  \lr{M\textsubscript{X}’(0) = E[X]}
\item
  \lr{M\textsubscript{X}’’(0) = E[X\textsuperscript{2}]}
\item
  \lr{Var(X) = M\textsubscript{X}’’(0) - [M\textsubscript{X}’(0)]\textsuperscript{2}}
\end{itemize}

\PSsubsection{ویژگی‌ها}

\begin{enumerate}
\def\labelenumi{\arabic{enumi}.}
\item
  \textbf{یگانگی:} اگر \lr{M\textsubscript{X}(t) = M\textsubscript{Y}(t)} برای تمام \lr{t} در همسایگی 0 برقرار باشد، آنگاه \lr{X} و \lr{Y} توزیع یکسانی دارند.
\item
  \textbf{جمع متغیرهای تصادفی مستقل:} اگر \lr{X} و \lr{Y} مستقل باشند:\PSbr{}\lr{M\textsubscript{X+Y}(t) = M\textsubscript{X}(t) \ensuremath{\cdot} M\textsubscript{Y}(t)}
\item
  \textbf{تبدیل خطی:} \lr{M\textsubscript{aX+b}(t) = e\textsuperscript{bt} \ensuremath{\cdot} M\textsubscript{X}(at)}
\end{enumerate}

\PSsubsection{\lr{MGF} توزیع‌های رایج}

{\def\LTcaptype{none} % do not increment counter
\begin{longtable}[c]{@{}cc@{}}
\toprule\noalign{}
\rowcolor{PSnavy}\PSth{توزیع} & \PSth{\lr{MGF} یعنی \lr{M\textsubscript{X}(t)}} \\
\midrule\noalign{}
\endhead
\bottomrule\noalign{}
\endlastfoot
\lr{Bernoulli(p)} & \lr{1 - p + pe\textsuperscript{t}} \\
\lr{Binomial(n,p)} & (\lr{1 - p + pe\textsuperscript{\lr{t})}n} \\
\lr{Poisson(\ensuremath{\lambda})} & \lr{exp(\ensuremath{\lambda}(e\textsuperscript{t} - 1))} \\
\lr{Exponential(\ensuremath{\lambda})} & \lr{\ensuremath{\lambda}/(\ensuremath{\lambda} - t), t \textless{} \ensuremath{\lambda}} \\
\lr{Gaussian(\ensuremath{\mu},\ensuremath{\sigma}\textsuperscript{2})} & \lr{exp(\ensuremath{\mu}t + \ensuremath{\sigma}\textsuperscript{2}t\textsuperscript{2}/2)} \\
\end{longtable}
}

\PSsubsection{کاربرد مهندسی: جمع سیگنال‌های مستقل}

اگر مؤلفه‌های سیگنال \lr{X\textsubscript{1} \textasciitilde{} N(\ensuremath{\mu}\textsubscript{1}, \ensuremath{\sigma}\textsubscript{1}\textsuperscript{2})} و \lr{X\textsubscript{2} \textasciitilde{} N(\ensuremath{\mu}\textsubscript{2}, \ensuremath{\sigma}\textsubscript{2}\textsuperscript{2})} مستقل باشند:

\lr{M\textsubscript{X\textsubscript{1}+X\textsubscript{2}}(t) = M\textsubscript{X\textsubscript{1}}(t) \ensuremath{\cdot} M\textsubscript{X\textsubscript{2}}(t)}\PSbr{}= \lr{exp(\ensuremath{\mu}\textsubscript{1}t + \ensuremath{\sigma}\textsubscript{1}\textsuperscript{2}t\textsuperscript{2}/2) \ensuremath{\cdot} exp(\ensuremath{\mu}\textsubscript{2}t + \ensuremath{\sigma}\textsubscript{2}\textsuperscript{2}t\textsuperscript{2}/2)}\PSbr{}= \lr{exp((\ensuremath{\mu}\textsubscript{1}+\ensuremath{\mu}\textsubscript{2})t + (\ensuremath{\sigma}\textsubscript{1}\textsuperscript{2}+\ensuremath{\sigma}\textsubscript{2}\textsuperscript{2})t\textsuperscript{2}/2)}

این همان \lr{MGF} توزیع \lr{N(\ensuremath{\mu}\textsubscript{1}+\ensuremath{\mu}\textsubscript{2}, \ensuremath{\sigma}\textsubscript{1}\textsuperscript{2}+\ensuremath{\sigma}\textsubscript{2}\textsuperscript{2})} است. بنابراین \lr{X\textsubscript{1}+X\textsubscript{2} \textasciitilde{} N(\ensuremath{\mu}\textsubscript{1}+\ensuremath{\mu}\textsubscript{2}, \ensuremath{\sigma}\textsubscript{1}\textsuperscript{2}+\ensuremath{\sigma}\textsubscript{2}\textsuperscript{2})}.

\PSsection{4.10}{کاربردهای مهندسی گشتاورها}

\PSsubsection{کاربرد 1: مشخصه‌سازی توان نویز}

یک سیگنال نویز \lr{N(t)} نمونه‌برداری می‌شود. نمونه‌ها میانگین \lr{\ensuremath{\mu}\textsubscript{N}} و واریانس \lr{\ensuremath{\sigma}\textsubscript{N}\textsuperscript{2}} دارند.

\begin{itemize}
\tightlist
\item
  \textbf{مؤلفه \lr{DC} (بایاس):} \lr{\ensuremath{\mu}\textsubscript{N}} \lr{—} برای نویز ایده‌آل باید صفر باشد
\item
  \textbf{توان نویز \lr{AC}:} \lr{\ensuremath{\sigma}\textsubscript{N}\textsuperscript{2}} \lr{—} تعیین‌کننده \lr{SNR}
\item
  \textbf{توان کل:} \lr{E[N\textsuperscript{2}] = \ensuremath{\mu}\textsubscript{N}\textsuperscript{2} + \ensuremath{\sigma}\textsubscript{N}\textsuperscript{2}}
\end{itemize}

\textbf{\lr{SNR} (نسبت سیگنال به نویز):}\PSdisplay{\PSmtext{SNR} = \frac{E[S^2]}{\sigma_N^2} = \frac{\PSmtext{Signal Power}}{\PSmtext{Noise Power}}}

\PSsubsection{کاربرد 2: خطای اندازه‌گیری}

یک حسگر مقدار واقعی \lr{\ensuremath{\theta}} را با خطای \lr{E} اندازه‌گیری می‌کند:

\begin{itemize}
\tightlist
\item
  اندازه‌گیری: \lr{X = \ensuremath{\theta} + E}
\item
  \lr{E[X] = \ensuremath{\theta} + E[E]} \lr{—} بایاس در اندازه‌گیری
\item
  \lr{Var(X) = Var(E)} \lr{—} دقت اندازه‌گیری
\end{itemize}

\textbf{صحّت:} \lr{\textbar{}}\lr{E[E]}\lr{\textbar{}} (میانگین چقدر به مقدار واقعی نزدیک است)\PSbr{}\textbf{دقت (تکرارپذیری):} \lr{\ensuremath{\sigma}\textsubscript{E}} (اندازه‌گیری‌ها چقدر تکرارپذیر هستند)

\PSsubsection{کاربرد 3: آماره‌های عمر قطعه}

قطعات عمر \lr{T \textasciitilde{} Exponential(\ensuremath{\lambda})} دارند:

\begin{itemize}
\tightlist
\item
  میانگین زمان تا خرابی: \lr{MTTF = E[T] = 1/\ensuremath{\lambda}}
\item
  انحراف معیار: \lr{\ensuremath{\sigma}\textsubscript{T} = 1/\ensuremath{\lambda} = MTTF}
\item
  برای نمایی: ضریب تغییرات \lr{CV = \ensuremath{\sigma}/\ensuremath{\mu} = 1} (همیشه!)
\end{itemize}

\PSsubsection{کاربرد 4: دامنه سیگنال مخابراتی}

سیگنال دریافتی \lr{R = A + N} که در آن \lr{A} = دامنه سیگنال و \lr{N \textasciitilde{} N(0, \ensuremath{\sigma}\textsuperscript{2})}:

\begin{itemize}
\tightlist
\item
  \lr{E[R] = A} (دامنه سیگنال به‌طور میانگین حفظ می‌شود)
\item
  \lr{Var(R) = \ensuremath{\sigma}\textsuperscript{2}} (توان نویز واریانس را تعیین می‌کند)
\item
  \lr{SNR = A\textsuperscript{2}/\ensuremath{\sigma}\textsuperscript{2}}
\end{itemize}

\PSsection{4.11}{مثال‌های \lr{MATLAB}}

\begin{PSbox}[unbreakable]{ex}{مثال 1: محاسبه گشتاورها از یک توزیع}

\tcbset{PScodemode/.style={unbreakable}}

\begin{Shaded}
\begin{Highlighting}[]
\CommentTok{\%\% Moments of Common Distributions}

\CommentTok{\% Exponential: failure time, lambda = 0.01 (per hour)}
\VariableTok{lambda} \OperatorTok{=} \FloatTok{0.01}\OperatorTok{;}
\VariableTok{beta} \OperatorTok{=} \FloatTok{1}\OperatorTok{/}\VariableTok{lambda}\OperatorTok{;}  \CommentTok{\% mean = 100 hours}

\VariableTok{fprintf}\NormalTok{(}\SpecialStringTok{\textquotesingle{}=== Exponential Distribution (λ = \%.3f) ===\textbackslash{}n\textquotesingle{}}\OperatorTok{,} \VariableTok{lambda}\NormalTok{)}\OperatorTok{;}
\VariableTok{fprintf}\NormalTok{(}\SpecialStringTok{\textquotesingle{}Mean (MTTF) = \%.1f hours\textbackslash{}n\textquotesingle{}}\OperatorTok{,} \VariableTok{beta}\NormalTok{)}\OperatorTok{;}
\VariableTok{fprintf}\NormalTok{(}\SpecialStringTok{\textquotesingle{}Variance = \%.1f hours²\textbackslash{}n\textquotesingle{}}\OperatorTok{,} \VariableTok{beta}\OperatorTok{\^{}}\FloatTok{2}\NormalTok{)}\OperatorTok{;}
\VariableTok{fprintf}\NormalTok{(}\SpecialStringTok{\textquotesingle{}Std Dev = \%.1f hours\textbackslash{}n\textquotesingle{}}\OperatorTok{,} \VariableTok{beta}\NormalTok{)}\OperatorTok{;}
\VariableTok{fprintf}\NormalTok{(}\SpecialStringTok{\textquotesingle{}Coefficient of Variation = \%.2f\textbackslash{}n\textquotesingle{}}\OperatorTok{,} \VariableTok{beta}\OperatorTok{/}\VariableTok{beta}\NormalTok{)}\OperatorTok{;}

\CommentTok{\% Gaussian: noise voltage, sigma = 0.1V}
\VariableTok{mu} \OperatorTok{=} \FloatTok{0}\OperatorTok{;} \VariableTok{sigma} \OperatorTok{=} \FloatTok{0.1}\OperatorTok{;}
\VariableTok{fprintf}\NormalTok{(}\SpecialStringTok{\textquotesingle{}\textbackslash{}n=== Gaussian Distribution (μ=\%.1f, σ=\%.2f V) ===\textbackslash{}n\textquotesingle{}}\OperatorTok{,} \VariableTok{mu}\OperatorTok{,} \VariableTok{sigma}\NormalTok{)}\OperatorTok{;}
\VariableTok{fprintf}\NormalTok{(}\SpecialStringTok{\textquotesingle{}Mean = \%.2f V\textbackslash{}n\textquotesingle{}}\OperatorTok{,} \VariableTok{mu}\NormalTok{)}\OperatorTok{;}
\VariableTok{fprintf}\NormalTok{(}\SpecialStringTok{\textquotesingle{}Variance (noise power) = \%.4f V²\textbackslash{}n\textquotesingle{}}\OperatorTok{,} \VariableTok{sigma}\OperatorTok{\^{}}\FloatTok{2}\NormalTok{)}\OperatorTok{;}
\VariableTok{fprintf}\NormalTok{(}\SpecialStringTok{\textquotesingle{}RMS noise = \%.3f V\textbackslash{}n\textquotesingle{}}\OperatorTok{,} \VariableTok{sigma}\NormalTok{)}\OperatorTok{;}
\VariableTok{fprintf}\NormalTok{(}\SpecialStringTok{\textquotesingle{}E[X²] (total power) = \%.4f V²\textbackslash{}n\textquotesingle{}}\OperatorTok{,} \VariableTok{mu}\OperatorTok{\^{}}\FloatTok{2} \OperatorTok{+} \VariableTok{sigma}\OperatorTok{\^{}}\FloatTok{2}\NormalTok{)}\OperatorTok{;}
\end{Highlighting}
\end{Shaded}

\end{PSbox}

\begin{PSbox}[unbreakable]{ex}{مثال 2: بررسی ویژگی‌های واریانس}

\tcbset{PScodemode/.style={unbreakable}}

\begin{Shaded}
\begin{Highlighting}[]
\CommentTok{\%\% Variance of Sum of Independent RVs}
\VariableTok{N} \OperatorTok{=} \FloatTok{100000}\OperatorTok{;}

\CommentTok{\% Two independent noise sources}
\VariableTok{sigma1} \OperatorTok{=} \FloatTok{0.1}\OperatorTok{;}   \CommentTok{\% RMS of noise 1}
\VariableTok{sigma2} \OperatorTok{=} \FloatTok{0.2}\OperatorTok{;}   \CommentTok{\% RMS of noise 2}

\VariableTok{X1} \OperatorTok{=} \VariableTok{sigma1} \OperatorTok{*} \VariableTok{randn}\NormalTok{(}\FloatTok{1}\OperatorTok{,} \VariableTok{N}\NormalTok{)}\OperatorTok{;}
\VariableTok{X2} \OperatorTok{=} \VariableTok{sigma2} \OperatorTok{*} \VariableTok{randn}\NormalTok{(}\FloatTok{1}\OperatorTok{,} \VariableTok{N}\NormalTok{)}\OperatorTok{;}
\VariableTok{Y} \OperatorTok{=} \VariableTok{X1} \OperatorTok{+} \VariableTok{X2}\OperatorTok{;}   \CommentTok{\% Combined noise}

\VariableTok{fprintf}\NormalTok{(}\SpecialStringTok{\textquotesingle{}Var(X1) = \%.4f (theory: \%.4f)\textbackslash{}n\textquotesingle{}}\OperatorTok{,} \VariableTok{var}\NormalTok{(}\VariableTok{X1}\NormalTok{)}\OperatorTok{,} \VariableTok{sigma1}\OperatorTok{\^{}}\FloatTok{2}\NormalTok{)}\OperatorTok{;}
\VariableTok{fprintf}\NormalTok{(}\SpecialStringTok{\textquotesingle{}Var(X2) = \%.4f (theory: \%.4f)\textbackslash{}n\textquotesingle{}}\OperatorTok{,} \VariableTok{var}\NormalTok{(}\VariableTok{X2}\NormalTok{)}\OperatorTok{,} \VariableTok{sigma2}\OperatorTok{\^{}}\FloatTok{2}\NormalTok{)}\OperatorTok{;}
\VariableTok{fprintf}\NormalTok{(}\SpecialStringTok{\textquotesingle{}Var(X1+X2) = \%.4f (theory: \%.4f)\textbackslash{}n\textquotesingle{}}\OperatorTok{,} \VariableTok{var}\NormalTok{(}\VariableTok{Y}\NormalTok{)}\OperatorTok{,} \VariableTok{sigma1}\OperatorTok{\^{}}\FloatTok{2} \OperatorTok{+} \VariableTok{sigma2}\OperatorTok{\^{}}\FloatTok{2}\NormalTok{)}\OperatorTok{;}
\VariableTok{fprintf}\NormalTok{(}\SpecialStringTok{\textquotesingle{}Std(X1+X2) = \%.4f (theory: \%.4f)\textbackslash{}n\textquotesingle{}}\OperatorTok{,} \VariableTok{std}\NormalTok{(}\VariableTok{Y}\NormalTok{)}\OperatorTok{,} \VariableTok{sqrt}\NormalTok{(}\VariableTok{sigma1}\OperatorTok{\^{}}\FloatTok{2}\OperatorTok{+}\VariableTok{sigma2}\OperatorTok{\^{}}\FloatTok{2}\NormalTok{))}\OperatorTok{;}
\VariableTok{fprintf}\NormalTok{(}\SpecialStringTok{\textquotesingle{}\textbackslash{}nNote: std devs do NOT add: \%.3f + \%.3f = \%.3f ≠ \%.3f\textbackslash{}n\textquotesingle{}}\OperatorTok{,} \OperatorTok{...}
    \VariableTok{sigma1}\OperatorTok{,} \VariableTok{sigma2}\OperatorTok{,} \VariableTok{sigma1}\OperatorTok{+}\VariableTok{sigma2}\OperatorTok{,} \VariableTok{sqrt}\NormalTok{(}\VariableTok{sigma1}\OperatorTok{\^{}}\FloatTok{2}\OperatorTok{+}\VariableTok{sigma2}\OperatorTok{\^{}}\FloatTok{2}\NormalTok{))}\OperatorTok{;}
\end{Highlighting}
\end{Shaded}

\end{PSbox}

\begin{PSbox}[unbreakable]{ex}{مثال 3: میانگین و واریانس یک متغیر تصادفی تبدیل‌شده}

\tcbset{PScodemode/.style={unbreakable}}

\begin{Shaded}
\begin{Highlighting}[]
\CommentTok{\%\% E[g(X)] using LOTUS: Power from Voltage}
\CommentTok{\% X = voltage \textasciitilde{} N(0, 0.5V)}
\CommentTok{\% P = X\^{}2 / R, with R = 50 ohms}

\VariableTok{sigma} \OperatorTok{=} \FloatTok{0.5}\OperatorTok{;}    \CommentTok{\% noise voltage RMS}
\VariableTok{R} \OperatorTok{=} \FloatTok{50}\OperatorTok{;}         \CommentTok{\% resistance}
\VariableTok{N} \OperatorTok{=} \FloatTok{100000}\OperatorTok{;}

\VariableTok{X} \OperatorTok{=} \VariableTok{sigma} \OperatorTok{*} \VariableTok{randn}\NormalTok{(}\FloatTok{1}\OperatorTok{,} \VariableTok{N}\NormalTok{)}\OperatorTok{;}
\VariableTok{P} \OperatorTok{=} \VariableTok{X}\OperatorTok{.\^{}}\FloatTok{2} \OperatorTok{/} \VariableTok{R}\OperatorTok{;}   \CommentTok{\% instantaneous power}

\CommentTok{\% Using LOTUS: E[X\^{}2/R] = E[X\^{}2]/R = sigma\^{}2/R}
\VariableTok{E\_P\_theory} \OperatorTok{=} \VariableTok{sigma}\OperatorTok{\^{}}\FloatTok{2} \OperatorTok{/} \VariableTok{R}\OperatorTok{;}
\VariableTok{E\_P\_sim} \OperatorTok{=} \VariableTok{mean}\NormalTok{(}\VariableTok{P}\NormalTok{)}\OperatorTok{;}

\VariableTok{fprintf}\NormalTok{(}\SpecialStringTok{\textquotesingle{}Mean power: Theory = \%.6f W, Sim = \%.6f W\textbackslash{}n\textquotesingle{}}\OperatorTok{,} \VariableTok{E\_P\_theory}\OperatorTok{,} \VariableTok{E\_P\_sim}\NormalTok{)}\OperatorTok{;}
\VariableTok{fprintf}\NormalTok{(}\SpecialStringTok{\textquotesingle{}Mean power = \%.3f mW\textbackslash{}n\textquotesingle{}}\OperatorTok{,} \VariableTok{E\_P\_theory} \OperatorTok{*} \FloatTok{1000}\NormalTok{)}\OperatorTok{;}
\end{Highlighting}
\end{Shaded}

\end{PSbox}

\begin{PSbox}[unbreakable]{ex}{مثال 4: گشتاورهای مراتب بالاتر و چولگی}

\tcbset{PScodemode/.style={unbreakable}}

\begin{Shaded}
\begin{Highlighting}[]
\CommentTok{\%\% Skewness and Kurtosis Comparison}
\VariableTok{N} \OperatorTok{=} \FloatTok{100000}\OperatorTok{;}

\CommentTok{\% Gaussian (symmetric)}
\VariableTok{X\_gauss} \OperatorTok{=} \VariableTok{randn}\NormalTok{(}\FloatTok{1}\OperatorTok{,} \VariableTok{N}\NormalTok{)}\OperatorTok{;}

\CommentTok{\% Exponential (right{-}skewed)}
\VariableTok{X\_exp} \OperatorTok{=} \VariableTok{exprnd}\NormalTok{(}\FloatTok{1}\OperatorTok{,} \FloatTok{1}\OperatorTok{,} \VariableTok{N}\NormalTok{)}\OperatorTok{;}

\CommentTok{\% Uniform (symmetric, light tails)}
\VariableTok{X\_unif} \OperatorTok{=} \VariableTok{rand}\NormalTok{(}\FloatTok{1}\OperatorTok{,} \VariableTok{N}\NormalTok{)}\OperatorTok{;}

\VariableTok{fprintf}\NormalTok{(}\SpecialStringTok{\textquotesingle{}Distribution    | Skewness | Kurtosis\textbackslash{}n\textquotesingle{}}\NormalTok{)}\OperatorTok{;}
\VariableTok{fprintf}\NormalTok{(}\SpecialStringTok{\textquotesingle{}{-}{-}{-}{-}{-}{-}{-}{-}{-}{-}{-}{-}{-}{-}{-}{-}+{-}{-}{-}{-}{-}{-}{-}{-}{-}{-}+{-}{-}{-}{-}{-}{-}{-}{-}{-}\textbackslash{}n\textquotesingle{}}\NormalTok{)}\OperatorTok{;}
\VariableTok{fprintf}\NormalTok{(}\SpecialStringTok{\textquotesingle{}Gaussian        | \%+.3f   | \%+.3f\textbackslash{}n\textquotesingle{}}\OperatorTok{,} \VariableTok{skewness}\NormalTok{(}\VariableTok{X\_gauss}\NormalTok{)}\OperatorTok{,} \VariableTok{kurtosis}\NormalTok{(}\VariableTok{X\_gauss}\NormalTok{)}\OperatorTok{{-}}\FloatTok{3}\NormalTok{)}\OperatorTok{;}
\VariableTok{fprintf}\NormalTok{(}\SpecialStringTok{\textquotesingle{}Exponential     | \%+.3f   | \%+.3f\textbackslash{}n\textquotesingle{}}\OperatorTok{,} \VariableTok{skewness}\NormalTok{(}\VariableTok{X\_exp}\NormalTok{)}\OperatorTok{,} \VariableTok{kurtosis}\NormalTok{(}\VariableTok{X\_exp}\NormalTok{)}\OperatorTok{{-}}\FloatTok{3}\NormalTok{)}\OperatorTok{;}
\VariableTok{fprintf}\NormalTok{(}\SpecialStringTok{\textquotesingle{}Uniform         | \%+.3f   | \%+.3f\textbackslash{}n\textquotesingle{}}\OperatorTok{,} \VariableTok{skewness}\NormalTok{(}\VariableTok{X\_unif}\NormalTok{)}\OperatorTok{,} \VariableTok{kurtosis}\NormalTok{(}\VariableTok{X\_unif}\NormalTok{)}\OperatorTok{{-}}\FloatTok{3}\NormalTok{)}\OperatorTok{;}
\VariableTok{fprintf}\NormalTok{(}\SpecialStringTok{\textquotesingle{}\textbackslash{}nTheoretical:\textbackslash{}n\textquotesingle{}}\NormalTok{)}\OperatorTok{;}
\VariableTok{fprintf}\NormalTok{(}\SpecialStringTok{\textquotesingle{}Gaussian:    skew=0, kurt=0\textbackslash{}n\textquotesingle{}}\NormalTok{)}\OperatorTok{;}
\VariableTok{fprintf}\NormalTok{(}\SpecialStringTok{\textquotesingle{}Exponential: skew=2, kurt=6\textbackslash{}n\textquotesingle{}}\NormalTok{)}\OperatorTok{;}
\VariableTok{fprintf}\NormalTok{(}\SpecialStringTok{\textquotesingle{}Uniform:     skew=0, kurt={-}1.2\textbackslash{}n\textquotesingle{}}\NormalTok{)}\OperatorTok{;}
\end{Highlighting}
\end{Shaded}

\end{PSbox}

\begin{PSbox}[unbreakable]{ex}{مثال 5: بررسی \lr{MGF}}

\tcbset{PScodemode/.style={unbreakable}}

\begin{Shaded}
\begin{Highlighting}[]
\CommentTok{\%\% Verify MGF: Compute moments by differentiation (numerical)}
\CommentTok{\% For X \textasciitilde{} Exponential(lambda)}
\VariableTok{lambda\_rate} \OperatorTok{=} \FloatTok{2}\OperatorTok{;}
\VariableTok{beta\_scale} \OperatorTok{=} \FloatTok{1}\OperatorTok{/}\VariableTok{lambda\_rate}\OperatorTok{;}

\CommentTok{\% MGF: M(t) = lambda/(lambda {-} t)}
\CommentTok{\% M\textquotesingle{}(0) = E[X], M\textquotesingle{}\textquotesingle{}(0) = E[X\^{}2]}

\VariableTok{t} \OperatorTok{=} \FloatTok{0}\OperatorTok{;}
\VariableTok{dt} \OperatorTok{=} \FloatTok{1e{-}6}\OperatorTok{;}

\CommentTok{\% Numerical derivatives}
\VariableTok{M} \OperatorTok{=} \OperatorTok{@}\NormalTok{(}\VariableTok{t}\NormalTok{) }\VariableTok{lambda\_rate} \OperatorTok{./}\NormalTok{ (}\VariableTok{lambda\_rate} \OperatorTok{{-}} \VariableTok{t}\NormalTok{)}\OperatorTok{;}
\VariableTok{M\_prime} \OperatorTok{=}\NormalTok{ (}\VariableTok{M}\NormalTok{(}\VariableTok{t}\OperatorTok{+}\VariableTok{dt}\NormalTok{) }\OperatorTok{{-}} \VariableTok{M}\NormalTok{(}\VariableTok{t}\OperatorTok{{-}}\VariableTok{dt}\NormalTok{)) }\OperatorTok{/}\NormalTok{ (}\FloatTok{2}\OperatorTok{*}\VariableTok{dt}\NormalTok{)}\OperatorTok{;}
\VariableTok{M\_double\_prime} \OperatorTok{=}\NormalTok{ (}\VariableTok{M}\NormalTok{(}\VariableTok{t}\OperatorTok{+}\VariableTok{dt}\NormalTok{) }\OperatorTok{{-}} \FloatTok{2}\OperatorTok{*}\VariableTok{M}\NormalTok{(}\VariableTok{t}\NormalTok{) }\OperatorTok{+} \VariableTok{M}\NormalTok{(}\VariableTok{t}\OperatorTok{{-}}\VariableTok{dt}\NormalTok{)) }\OperatorTok{/} \VariableTok{dt}\OperatorTok{\^{}}\FloatTok{2}\OperatorTok{;}

\VariableTok{fprintf}\NormalTok{(}\SpecialStringTok{\textquotesingle{}From MGF derivatives:\textbackslash{}n\textquotesingle{}}\NormalTok{)}\OperatorTok{;}
\VariableTok{fprintf}\NormalTok{(}\SpecialStringTok{\textquotesingle{}  E[X] = M\textquotesingle{}\textquotesingle{}(0) = \%.4f (theory: \%.4f)\textbackslash{}n\textquotesingle{}}\OperatorTok{,} \VariableTok{M\_prime}\OperatorTok{,} \FloatTok{1}\OperatorTok{/}\VariableTok{lambda\_rate}\NormalTok{)}\OperatorTok{;}
\VariableTok{fprintf}\NormalTok{(}\SpecialStringTok{\textquotesingle{}  E[X²] = M\textquotesingle{}\textquotesingle{}\textquotesingle{}\textquotesingle{}(0) = \%.4f (theory: \%.4f)\textbackslash{}n\textquotesingle{}}\OperatorTok{,} \VariableTok{M\_double\_prime}\OperatorTok{,} \FloatTok{2}\OperatorTok{/}\VariableTok{lambda\_rate}\OperatorTok{\^{}}\FloatTok{2}\NormalTok{)}\OperatorTok{;}
\VariableTok{fprintf}\NormalTok{(}\SpecialStringTok{\textquotesingle{}  Var(X) = E[X²] {-} (E[X])² = \%.4f (theory: \%.4f)\textbackslash{}n\textquotesingle{}}\OperatorTok{,} \OperatorTok{...}
    \VariableTok{M\_double\_prime} \OperatorTok{{-}} \VariableTok{M\_prime}\OperatorTok{\^{}}\FloatTok{2}\OperatorTok{,} \FloatTok{1}\OperatorTok{/}\VariableTok{lambda\_rate}\OperatorTok{\^{}}\FloatTok{2}\NormalTok{)}\OperatorTok{;}
\end{Highlighting}
\end{Shaded}

\end{PSbox}

\PSsection{4.12}{مسایل تمرینی}

\begin{PSbox}[unbreakable]{prob}{مسیله 1: محاسبه پایه}

یک متغیر تصادفی گسسته \lr{X} دارای \lr{PMF} زیر است: \lr{P(X=-1) = 0.2}، \lr{P(X=0) = 0.5}، \lr{P(X=1) = 0.2}، \lr{P(X=2) = 0.1}.

\begin{enumerate}
\def\labelenumi{(\alph{enumi})}
\tightlist
\item
  \lr{E[X]} را بیابید.
\item
  \lr{E[X\textsuperscript{2}]} را بیابید.
\item
  \lr{Var(X)} را بیابید.
\item
  \lr{E[3X + 2]} را بیابید.
\item
  \lr{Var(3X + 2)} را بیابید.
\end{enumerate}

\PSsolution{} 

\begin{enumerate}
\def\labelenumi{(\alph{enumi})}
\item
  \lr{E[X] = (-1)(0.2) + (0)(0.5) + (1)(0.2) + (2)(0.1) = -0.2 + 0 + 0.2 + 0.2 = 0.2}
\item
  \lr{E[X\textsuperscript{2}] = (1)(0.2) + (0)(0.5) + (1)(0.2) + (4)(0.1) = 0.2 + 0 + 0.2 + 0.4 = 0.8}
\item
  \lr{Var(X) = E[X\textsuperscript{2}] - (E[X])\textsuperscript{2} = 0.8 - 0.04 = 0.76}
\item
  \lr{E[3X + 2] = 3E[X] + 2 = 3(0.2) + 2 = 2.6}
\item
  \lr{Var(3X + 2) = 9\ensuremath{\cdot}Var(X) = 9(0.76) = 6.84}
\end{enumerate}

\end{PSbox}

\begin{PSbox}[unbreakable]{prob}{مسیله 2: سیگنال و نویز}

یک سیگنال دریافتی \lr{R = 2.5 + N} که در آن \lr{N \textasciitilde{} N(0, 0.04)}.

\begin{enumerate}
\def\labelenumi{(\alph{enumi})}
\tightlist
\item
  \lr{E[R]} و \lr{Var(R)} را بیابید.
\item
  \lr{SNR} را بر حسب \lr{dB} بیابید.
\item
  \lr{P(R \textless{} 0)} را بیابید \lr{—} احتمال این‌که سیگنال منفی شود.
\end{enumerate}

\PSsolution{} 

\begin{enumerate}
\def\labelenumi{(\alph{enumi})}
\item
  \lr{E[R] = 2.5}، \lr{Var(R) = Var(N) = 0.04}، \lr{\ensuremath{\sigma}\textsubscript{R} = 0.2}
\item
  توان سیگنال = \lr{2.5\textsuperscript{2} = 6.25}، توان نویز = \PSnum{0.04}\PSbr{}\lr{SNR = 6.25/0.04 = 156.25 \ensuremath{\to} SNR\textsubscript{dB} = 10\ensuremath{\cdot}log\textsubscript{10}(156.25) = 21.94 dB}
\item
  \lr{P(R \textless{} 0) = P(N \textless{} -2.5) = \ensuremath{\Phi}(-2.5/0.2) = \ensuremath{\Phi}(-12.5) \ensuremath{\approx} 0} (بسیار نامحتمل)
\end{enumerate}

\end{PSbox}

\begin{PSbox}[unbreakable]{prob}{مسیله 3: میانگین‌گیری از اندازه‌گیری‌ها}

یک حسگر \lr{n} اندازه‌گیری مستقل از یک ولتاژ ثابت \lr{V = 3.3V} انجام می‌دهد.\PSbr{}هر اندازه‌گیری: \lr{X\textsubscript{i} = 3.3 + N\textsubscript{i}} که در آن \lr{N\textsubscript{i} \textasciitilde{} N(0, 0.01)}.

\begin{enumerate}
\def\labelenumi{(\alph{enumi})}
\tightlist
\item
  برای \lr{n = 1} مقدار \lr{E[\ensuremath{\bar{X}}]} و \lr{Var(\ensuremath{\bar{X}})} چقدر است؟
\item
  چه تعداد اندازه‌گیری \lr{n} لازم است تا \lr{\ensuremath{\sigma}\textsubscript{X}̄ \textless{} 0.01V} شود؟
\item
  با \lr{n = 100}، مقدار \lr{P(\textbar{}\ensuremath{\bar{X}} - 3.3\textbar{} \textgreater{} 0.02)} چقدر است؟
\end{enumerate}

\PSsolution{} 

\begin{enumerate}
\def\labelenumi{(\alph{enumi})}
\item
  \lr{E[\ensuremath{\bar{X}}] = 3.3}، \lr{Var(\ensuremath{\bar{X}}) = 0.01/1 = 0.01}، \lr{\ensuremath{\sigma} = 0.1V}
\item
  \lr{\ensuremath{\sigma}\textsubscript{X}̄ = \ensuremath{\sigma}/\ensuremath{\surd}n \textless{} 0.01 \ensuremath{\to} \ensuremath{\surd}n \textgreater{} 10 \ensuremath{\to} n \textgreater{} 100 \ensuremath{\to}} به \lr{n \ensuremath{\ge} 101} نیاز است
\item
  \lr{\ensuremath{\sigma}\textsubscript{X}̄ = 0.1/\ensuremath{\surd}100 = 0.01V}\PSbr{}\lr{P(\textbar{}\ensuremath{\bar{X}} - 3.3\textbar{} \textgreater{} 0.02) = P(\textbar{}Z\textbar{} \textgreater{} 2) = 2(1-\ensuremath{\Phi}(2)) = 0.0455}
\end{enumerate}

\end{PSbox}

\begin{PSbox}[unbreakable]{prob}{مسیله 4: عمر قطعه}

قطعات عمر \lr{T \textasciitilde{} Exp(\ensuremath{\lambda} = 0.002 \rl{در هر ساعت})} دارند.

\begin{enumerate}
\def\labelenumi{(\alph{enumi})}
\tightlist
\item
  \lr{MTTF}، \lr{\ensuremath{\sigma}\textsubscript{T}} و ضریب تغییرات را بیابید.
\item
  \lr{P(T \textgreater{} 1000 \rl{ساعت})} را بیابید.
\item
  \lr{E[T\textsuperscript{2}]} را بیابید و تفسیر کنید.
\end{enumerate}

\PSsolution{} 

\begin{enumerate}
\def\labelenumi{(\alph{enumi})}
\item
  \lr{MTTF = 1/\ensuremath{\lambda} = 500} ساعت، \lr{\ensuremath{\sigma}\textsubscript{T} = 500} ساعت، \lr{CV = \ensuremath{\sigma}/\ensuremath{\mu} = 1}
\item
  \lr{P(T \textgreater{} 1000) = e\textsuperscript{-0.002\ensuremath{\times}1000} = e\textsuperscript{-2} = 0.1353}
\item
  \lr{E[T\textsuperscript{2}] = Var(T) + (E[T])\textsuperscript{2} = 500\textsuperscript{2} + 500\textsuperscript{2} = 500000 hours\textsuperscript{2}}\PSbr{}\lr{\ensuremath{\surd}E[T\textsuperscript{2}] = 707} ساعت (عمر مؤثر)
\end{enumerate}

\end{PSbox}

\begin{PSbox}[unbreakable]{prob}{مسیله 5: کاربرد \lr{MGF}}

\lr{X} دارای \lr{MGF} زیر است: \lr{M\textsubscript{X}(t) = (1/3)e\textsuperscript{t} + (2/3)e\textsuperscript{2t}}.

\begin{enumerate}
\def\labelenumi{(\alph{enumi})}
\tightlist
\item
  \lr{X} چه نوع متغیر تصادفی است؟ \lr{PMF} آن را بیابید.
\item
  با استفاده از \lr{MGF} مقدار \lr{E[X]} را بیابید.
\item
  با استفاده از \lr{MGF} مقدار \lr{Var(X)} را بیابید.
\end{enumerate}

\PSsolution{} 

\begin{enumerate}
\def\labelenumi{(\alph{enumi})}
\item
  \lr{X} گسسته است. با مقایسه با \lr{\ensuremath{\Sigma} p\textsubscript{X}(x)e\textsuperscript{tx}}:\PSbr{}\lr{P(X=1) = 1/3}، \lr{P(X=2) = 2/3}
\item
  \lr{M’(t) = (1/3)e\textsuperscript{t} + (4/3)e\textsuperscript{2t}}\PSbr{}\lr{E[X] = M’(0) = 1/3 + 4/3 = 5/3}
\item
  \lr{M’«\lr{(t) = (1/3)e\textsuperscript{t} + (8/3)e\textsuperscript{2t}}\PSbr{}\lr{E[X\textsuperscript{2}] = M}»’(0) = 1/3 + 8/3 = 3}\PSbr{}\lr{Var(X) = 3 - (5/3)\textsuperscript{2} = 3 - 25/9 = 2/9}
\end{enumerate}

\end{PSbox}

\PSsection{4.13}{نکات کلیدی}

\begin{enumerate}
\def\labelenumi{\arabic{enumi}.}
\tightlist
\item
  \textbf{\lr{E[X]} = میانگین درازمدت} \lr{—} مرکز توزیع، مؤلفه \lr{DC}
\item
  \textbf{\lr{Var(X)} = پراکندگی/توان} \lr{—} برای سیگنال‌های با میانگین صفر، واریانس همان توان است
\item
  \textbf{\lr{\ensuremath{\sigma}} همان یکای \lr{X} را دارد} \lr{—} مستقیماً به‌عنوان مقدار مؤثر \lr{(RMS)} قابل تفسیر است
\item
  \textbf{خطی‌بودن امید ریاضی} همواره برقرار است؛ جمع واریانس‌ها نیازمند استقلال است
\item
  \textbf{\lr{MGF} به‌طور یکتا} یک توزیع را مشخص می‌کند و تمام گشتاورها را تولید می‌کند
\item
  \textbf{در مهندسی:} میانگین \lr{\ensuremath{\to}} سطح سیگنال، واریانس \lr{\ensuremath{\to}} توان نویز، انحراف معیار \lr{\ensuremath{\to}} نویز مؤثر
\end{enumerate}

\emph{ماژول بعدی: {ماژول 5 \lr{—} توابع مهم متغیر تصادفی}}\PSbr{}
